% BEGIN REV09_TRANSPORTE_POLAR
\subsection{Transporte polar de la incidencia y conservación de los lectores}
\label{corpus9:iv:polar}

La realización finita del complejo orientado de incidencia
recibe la frontera y la orientación construidas en el estado
APP--TRIT--TPK. Una vez fijados los productos escalares de
esa realización, sean \(\mathcal H_0,\mathcal H_1\) sus
espacios finitos y
\[
 d:\mathcal H_0\longrightarrow\mathcal H_1,\qquad
 \Delta_0=d^*d,\qquad \Delta_1=dd^*.
\]
Los lectores de incidencia y retorno son
\[
 \rho_{\mathrm{int}}(v)=(dv,d^*dv),\qquad
 \rho_{\mathrm{ref}}(a)=(d^*a,dd^*a).
\]
Los símbolos \(\Delta_0,\Delta_1\) designan aquí los
Laplacianos de esta realización hilbertiana. La orientación,
las métricas y el dominio de \(d\) forman parte de sus datos.

\begin{proposition}[Equivalencia unitaria de los sectores positivos]
\label{corpus9:iv:polar-transporte}
Sean
\[
 \mathcal H_0^+=(\ker d)^\perp,\qquad
 \mathcal H_1^+=(\ker d^*)^\perp=\operatorname{Ran}d.
\]
La aplicación
\begin{equation}
 \mathcal U=d(\Delta_0|_{\mathcal H_0^+})^{-1/2}
    :\mathcal H_0^+\longrightarrow\mathcal H_1^+
 \label{corpus9:iv:polar-operador}
\end{equation}
es unitaria y satisface, en esos sectores,
\begin{equation}
 \mathcal U\Delta_0=\Delta_1\mathcal U,\qquad
 \mathcal U f(\Delta_0)=f(\Delta_1)\mathcal U
 \label{corpus9:iv:polar-funcional}
\end{equation}
para toda función \(f\) definida sobre el espectro positivo
finito común.
\end{proposition}
\begin{proof}
La descomposición singular reducida es \(d=V\Sigma W^*\),
donde las columnas de \(W\) y \(V\) son bases ortonormales
de \(\mathcal H_0^+\) y \(\mathcal H_1^+\), y \(\Sigma\)
es diagonal con valores estrictamente positivos. Por tanto,
\[
 \Delta_0|_{\mathcal H_0^+}=W\Sigma^2W^*,\qquad
 \Delta_1|_{\mathcal H_1^+}=V\Sigma^2V^*,\qquad
 \mathcal U=VW^*.
\]
Se sigue \(\mathcal U^*\mathcal U=I_{\mathcal H_0^+}\) y
\(\mathcal U\mathcal U^*=I_{\mathcal H_1^+}\). Evaluar \(f\)
en la misma diagonal \(\Sigma^2\) prueba
\eqref{corpus9:iv:polar-funcional}. Si \(d=0\), ambos
sectores positivos son nulos y las identidades tienen su
interpretación única entre espacios de dimensión cero.
\end{proof}

Para \(v\in\mathcal H_0^+\) y \(h,t\geq0\), se conservan
la energía cuadrática, el resolvente y el semigrupo:
\begin{align}
 \langle\mathcal Uv,\Delta_1\mathcal Uv\rangle
    &=\langle v,\Delta_0v\rangle,\\
 \mathcal U(I+h\Delta_0)^{-1}
    &=(I+h\Delta_1)^{-1}\mathcal U,\\
 \mathcal Ue^{-t\Delta_0}
    &=e^{-t\Delta_1}\mathcal U.
 \label{corpus9:iv:polar-lectores}
\end{align}
La primera identidad se deduce de la unitariedad y del
entrelazamiento. Las otras dos son casos de
\eqref{corpus9:iv:polar-funcional}, con
\(f(x)=(1+hx)^{-1}\) y \(f(x)=e^{-tx}\).

\subsubsection{Realización sobre el ciclo trítico}
En las bases ortonormales de vértices y aristas orientadas
del ciclo de tres elementos,
\[
 d=
 \begin{pmatrix}
 -1&1&0\\
 0&-1&1\\
 1&0&-1
 \end{pmatrix}.
\]
La multiplicación de matrices da
\[
 \Delta_0=\Delta_1=
 \begin{pmatrix}
 2&-1&-1\\
 -1&2&-1\\
 -1&-1&2
 \end{pmatrix}
 =3I-J_3=2I-C-C^{-1},
\]
donde \(J_3\) tiene todas sus entradas iguales a uno y \(C\)
es la permutación cíclica. El vector constante genera el
núcleo; sobre su complemento ortogonal, \(J_3=0\).
Así, el espectro es \(\{0,3,3\}\) y
\(\mathcal U=d/\sqrt3\) en el sector positivo.
Los dos lectores comparten este espectro por proceder de la
misma incidencia, conservando sus espacios de vértices y
aristas.

\subsubsection{Orientación y compatibilidad de refinamientos}
La sustitución \(d\mapsto-d\) conserva \(\Delta_0,\Delta_1\)
y cambia \(\mathcal U\) por \(-\mathcal U\). Más
generalmente, para una inversión de orientaciones de aristas
representada por una matriz diagonal unitaria \(E\),
\[
 d'=Ed,\qquad \Delta_0'=\Delta_0,\qquad
 \Delta_1'=E\Delta_1E^*,\qquad \mathcal U'=E\mathcal U.
\]
Estas igualdades se obtienen usando \(E^*E=I\) en las
definiciones. La orientación queda así conservada por el
transportador, incluso cuando las lecturas espectrales
coinciden.

En una colección de refinamientos, los mapas
\(J_m^{(0)}:\mathcal H_{0,m}^+\to\mathcal H_{0,m+1}^+\)
y
\(J_m^{(1)}:\mathcal H_{1,m}^+\to\mathcal H_{1,m+1}^+\)
deben conservar los dominios y normas y satisfacer
\[
 \mathcal U_{m+1}J_m^{(0)}
    =J_m^{(1)}\mathcal U_m
\]
para transportar esta identificación a través de niveles.
La proposición prueba la equivalencia finita de los sectores
positivos; el cuadrado anterior expresa la condición adicional
del refinamiento. Los núcleos, las restricciones no lineales
y los pasos al límite conservan sus demostraciones propias.
% END REV09_TRANSPORTE_POLAR
